\section{Performance \& Resource Analysis}
\label{sec:performance_analysis}

In the PIN cracking scenario, the most significant performance bottleneck is the system latency incurred each time authentication is submitted in the form of blocked I/O inside the \texttt{submit\_pin} function. In order to maximize the efficiency, we optimize three aspects of the system --- algorithm architecture, os scheduling and underlying hardware interaction.

In terms of algorithm architecture, the program completely abandoned the resource-consuming naive brute-force search paradigm when dealing with the cracking task of arbitrary length sequences. The time complexity of brute-force search is as high as $O(D^N)$ (where \texttt{N} is the length of the sequence and \texttt{D} is the optional number base), which will lead to state space explosion in the case of long sequence and large base, extremely waste of CPU cycles and memory bandwidth. For this purpose, we use a constructive combinatorial search algorithm based on backtracking. It implements   a pruning by using the baseline Hamming distance returned by the initial reference sequence as an absolute constraint. The algorithm only deduces dynamically in memory and generates valid candidate sequences that strictly satisfy the Hamming distance characteristics of the target. The time complexity is dramatically compressed to $O(C(N,K) \times (D-1)^K)$ (where \texttt{K} is the known Hamming distance).

As for OS scheduling, exceptional control flow is introduced to optimize the resource consumption during I/O polling. The CPU-intensive timestamps are not used during the processing of button input to achieve accurate timeout detection with numeric separations. Instead, the system calls the interval timer \texttt{setitimer} of the Linux kernel and binds the \texttt{SIGALRM} asynchronous signal handler function, which enables the Raspberry Pi to effectively release resources to process other system-level tasks during the idle period of waiting for the user's key release or timeout, which greatly reduces the power consumption and unnecessary processor usage caused by single thread polling.

In terms of underlying hardware interaction and resource management, the program directly controls the physical hardware to minimize overhead. As will be detailed in the subsequent sections, implementing the high-frequency Hamming distance calculation in ARM assembly eliminates redundant instructions from high-level compilers. For GPIO control, mapping physical addresses via \texttt{mmap} and using inline assembly help to bypass library abstractions, ensuring microsecond I/O latency. Regarding memory, we avoid large static arrays; all sequence buffers are dynamically allocated based on the runtime \texttt{seqlen}. Task 5's recursive search bounds its maximum call stack depth to the target Hamming distance, which is typically much smaller than the total sequence length. This control over recursion, combined with on-demand memory allocation, effectively prevents stack overflows and maintains a minimal memory consumption when cracking extremely long sequences.
